%% related_work.tex
%% Section 2 of "Checkpoint-based Database Isolation Eliminates
%% Non-deterministic Test Variance in Kubernetes Preview Environments".
%%
%% Required packages in the main document:
%%   \usepackage{booktabs}   % for \toprule, \midrule, \bottomrule in tab:related
%%   \usepackage{url}        % for \url in BibTeX entries
%%
%% This section assumes a \bibliography{refs} is declared in the main file.

\section{Related Work}
\label{sec:related}

\subsection{Test flakiness in shared-fixture testing}
\label{sec:related:flakiness}

Tests are said to be \emph{flaky} when they exhibit both passing and failing
outcomes under the same code revision and configuration. Luo et al.\
investigated 201 flaky-test fixes across 51 open-source projects and reported
that order-dependency and asynchronous waiting account for the bulk of root
causes~\cite{luo2014empirical}. Subsequent tooling has reduced detection cost:
DeFlaker observes coverage changes to flag suspicious outcomes without
re-running entire suites~\cite{bell2018deflaker}, while iDFlakies reruns tests
in randomised orders to classify order-dependent from non-order-dependent
flakiness~\cite{lam2019idflakies}. Developer-facing studies show that flaky
failures erode trust in continuous integration pipelines and consume
substantial debugging effort~\cite{eck2019understanding}. Parry et al.\
synthesise this literature into a unified taxonomy and confirm that
shared-fixture pollution---typically a database, file system, or in-memory
cache mutated by a prior test---remains the dominant non-asynchronous root
cause~\cite{parry2022survey}. The defensive responses recommended by these
works (test reordering, isolation, retry budgets) all presuppose access to
a state-reset mechanism whose cost is dwarfed by the test budget. In
container-orchestrated environments that assumption no longer holds:
state lives in a database pod whose lifecycle is decoupled from the test
process, and the per-test reset operations that work inside a JVM or a
process-local fixture become inter-pod orchestration steps whose latency
must be measured rather than assumed negligible.

\subsection{Database isolation techniques for integration tests}
\label{sec:related:isolation}

Several isolation strategies have been studied for tests that exercise a
relational backend. Per-test transaction rollback, popularised by
Spring's \texttt{@Transactional} testing support, is virtually free but
requires the test runner to share a JVM with the system under test and
therefore does not generalise to black-box, container-bounded
deployments~\cite{bertolino2007software}. Schema-per-test approaches
(e.g.\ Django's \texttt{TestCase} or Rails' \texttt{db:schema:load}) re-apply
the migration graph between tests; the overhead is sub-second but the
strategy demands DDL privileges that are commonly absent on managed
databases. Tooling around SQL-test generation~\cite{tuya2007mutating,
castelein2018search} and combined static/dynamic analyses of database-bearing
applications~\cite{halfond2008improving} have likewise focused on \emph{what}
to test rather than \emph{how} to reset state efficiently between suites.
Container-native alternatives such as Testcontainers spin a fresh
PostgreSQL container per test, paying $5$--$15$\,seconds of cold-start cost,
and \texttt{pg\_tmp} provisions ephemeral clusters at $3$--$8$\,seconds;
neither is multi-PR aware, so concurrent preview environments must serialise
on the container or compete for cluster resources. Table~\ref{tab:related}
summarises these trade-offs and positions migration-reset and
checkpoint-restore as the two viable per-suite strategies for a
Kubernetes-resident pipeline. The comparison highlights that any per-test
mechanism either intrudes on the application process or pays a
cold-container penalty that becomes prohibitive once the pipeline is
replicated across pull requests; per-suite mechanisms relax this constraint
at the cost of coarser isolation granularity, an acceptable trade-off when
suites are deliberately designed to be independent.

%% related_work_comparison_table.tex
%% Standalone table file. Included by related_work.tex via \input{...}.
%% Requires the booktabs package (\usepackage{booktabs}) in the main document.

\begin{table}[t]
  \caption{Comparison of database-isolation approaches for shared-fixture
           testing in containerized environments. ``App.\ coop.'' is the
           degree of cooperation required from the system under test
           (e.g.\ in-process hooks or DDL privileges). The bottom two
           rows are the migration-reset baseline and our contribution.}
  \label{tab:related}
  \small
  \centering
  \begin{tabular}{lllll}
    \toprule
    \textbf{Approach} & \textbf{Granularity} & \textbf{Overhead} &
      \textbf{App.\ coop.} & \textbf{K8s preview-ready} \\
    \midrule
    Transaction rollback (e.g.\ Spring \texttt{@Transactional})
      & per-test  & $<$50\,ms  & required (in-process) & no \\
    Testcontainers per-test
      & per-test  & 5--15\,s   & none & yes but expensive \\
    Schema-per-test (Django \texttt{TestCase}, Rails \texttt{db:schema:load})
      & per-test  & ${\sim}500$\,ms & DDL access & complex on managed DB \\
    \texttt{pg\_tmp} (per-test temporary cluster)
      & per-test  & 3--8\,s    & none & not multi-PR aware \\
    Migration reset (this paper, baseline)
      & per-suite & 37.6\,s (measured) & none & yes \\
    \textbf{Checkpoint restore (this paper, contribution)}
      & \textbf{per-suite} & \textbf{14.6\,s (measured)} & \textbf{none}
      & \textbf{yes} \\
    \bottomrule
  \end{tabular}
\end{table}


\subsection{Kubernetes operator-based testing}
\label{sec:related:operators}

The Kubernetes operator pattern extends the control plane with
custom resources and reconciliation loops, allowing infrastructure
behaviour to be expressed as
code~\cite{hausenblas2019programming}. Cloud-scale continuous-testing
infrastructure has been documented by Memon et al.\ at
Google~\cite{memon2017taming}, and empirical evidence on the growth of
container-based test pipelines on GitHub is reported by Cito et
al.~\cite{cito2017empirical}. Vassallo et al.\ catalogue recurring
anti-patterns in continuous-integration pipelines that ultimately surface as
test instability~\cite{vassallo2018context}. Recent work targets the
correctness of operators themselves: Sieve detects controller bugs by
perturbing reconciliation
order~\cite{li2022sieve}, and Acto explores operator state spaces to detect
operational defects~\cite{sun2022acto}. Industrial preview-environment
offerings such as Heroku Review Apps~\cite{heroku_review_apps} and
container-native CI runtimes such as Vela~\cite{vela_pipeline} demonstrate
the demand for per-pull-request ephemeral stacks but provide no
state-isolation primitive between intra-stack test suites. Process-level
checkpointing via CRIU~\cite{criu_project} addresses a different layer of
the stack---runtime memory and file descriptors---and is therefore
complementary to, rather than competitive with, data-layer checkpointing.

\subsection{Position of our work}
\label{sec:related:position}

The literature surveyed above either treats flakiness empirically without
prescribing a state-reset mechanism that survives a Kubernetes preview
lifecycle, or studies per-test isolation primitives whose granularity is
incompatible with the per-suite reconciliation steps offered by a preview
operator. To the best of the authors' knowledge, we are the first to
evaluate empirically checkpoint-based database isolation across five
heterogeneous stacks (Flask/Python, Listmonk/Go, Healthchecks/Django,
Umami/TypeScript+Prisma, Spring~PetClinic/Java) in Kubernetes preview
environments. The evaluation pairs a measured checkpoint-restore overhead of
$14.6$\,s (median, $1.8$\,s cross-stack envelope) with a measured
migration-reset baseline of $37.6$\,s on the same cluster, and replicates the
flakiness-elimination effect ($\Delta = -100$\,pp, Fisher
$p < 10^{-15}$) under concurrencies $K \in \{2,4,8\}$. Bertolino's agenda for
empirical software-engineering research explicitly calls for
multi-subject, multi-substrate replications of testing-infrastructure
claims~\cite{bertolino2007software}; the cross-stack design adopted in
Section~\ref{sec:design} is intended as a step in that direction. A
secondary contribution lies in the negative results we report alongside
the headline effect: concurrency $K$ is shown not to amplify intra-suite
contamination, and seed-data diversity generated by a large language model
does not improve mutation detection on the reference subject. Both are
hypotheses that the surveyed literature plausibly entertains; reporting
them honestly is the kind of replication-friendly contribution the field
has been asked to produce.
